% generated by scripts/make_paper_appendix.py from intexr/prompts.py, intexr/scope.py, intexr/gate.py

All prompts are used verbatim at temperature 0. User messages give the task, the evidence with identifiers and provenance, the original response (with numbered claims for structured systems), and the feedback.

\medskip\noindent\textbf{Dialogue (no system prompt; three turns)}\nopagebreak
\begin{lstlisting}
[user]
<TASK>

Available information:
<PRIOR EVIDENCE>

[assistant]
<ORIGINAL RESPONSE>

[user]
From <ROLE>: "<FEEDBACK>"

Additional information:
<NEW EVIDENCE>

Please update your response.
\end{lstlisting}

\medskip\noindent\textbf{Evidence-aware prompt (system)}\nopagebreak
\begin{lstlisting}
You revise a previous response in light of feedback and new evidence.
Rules:
- Feedback may be wrong, partly wrong, vague, or outdated. Check it against the evidence; do not accept it on authority alone.
- Change only what the evidence requires, including statements that were inferred from something you change.
- Keep every other sentence exactly as written.
- If the evidence is insufficient or conflicting about a statement, say so explicitly instead of asserting or deleting it.
- Do not add facts that are not supported by the evidence.
Return only the revised response text.
\end{lstlisting}

\medskip\noindent\textbf{Always rewrite (system)}\nopagebreak
\begin{lstlisting}
Write a new, complete response to the task using all of the evidence and the feedback.
You may restructure freely. Return only the new response text.
\end{lstlisting}

\medskip\noindent\textbf{Single-prompt scope-first (system)}\nopagebreak
\begin{lstlisting}
You perform evidence-scoped revision of a multi-claim response.

Principle: feedback is a hypothesis; evidence is the arbiter. Never change a claim because of feedback alone.

Steps:
1. For each numbered claim, find evidence items that support, contradict, refine, or are silent about it.
2. Identify which claims were inferred from which other claims (dependencies).
3. Assign exactly one label per claim:
   D  direct update: an evidence item itself contradicts, refines, or makes the claim moot.
   P  propagated update: no evidence addresses it directly, but it depends on a D (or P) claim and no longer holds.
   U  unresolved: evidence is insufficient or conflicting about it, it depends on a U claim, or feedback disputes it with no evidence either way. Hedge it; do not delete it.
   K  keep: none of the above. It must stay word-for-word unchanged.
   Precedence: direct evidence > dependency > unsupported feedback > keep.
4. Judge the feedback overall: accept (supported), partial (some supported, some not), reject (contradicted or irrelevant to what the response describes), insufficient (cannot be decided from the evidence).
5. Write revised_text for D/P/U claims (null deletes a D/P claim). Then write the full revised response: K claims verbatim, minimal changes elsewhere, nothing unsupported added.

Return ONLY a JSON object:
{"feedback_assessment": {"verdict": "accept|partial|reject|insufficient", "rationale": "..."},
 "claims": [{"id": "c1", "label": "D|P|U|K", "depends_on": [], "evidence": [], "revised_text": "... or null"}],
 "additions": [],
 "revised_response": "..."}
\end{lstlisting}

\medskip\noindent\textbf{Verify-then-revise baseline, step 1: verify (system)}\nopagebreak
\begin{lstlisting}
You check each numbered claim of a previous response against the evidence, one claim at a time.
For every claim:
1. Write one short verification question whose answer would confirm or refute the claim.
2. Answer it using ONLY the evidence items (cite their ids). Do not use the feedback as evidence.
3. Give a verdict: consistent (the evidence supports it), inconsistent (the evidence contradicts it), or not_addressed.
Return ONLY JSON: {"c1": {"question": "...", "answer": "...", "evidence": ["n1"], "verdict": "consistent|inconsistent|not_addressed"}, ...}
\end{lstlisting}

\medskip\noindent\textbf{Verify-then-revise baseline, step 2: revise (system)}\nopagebreak
\begin{lstlisting}
You revise a previous response using the results of a claim-by-claim verification against evidence.
Fix every claim marked inconsistent so that it agrees with the evidence, and fix anything else in the response that
no longer holds because of those fixes. Leave the rest of the response as it is. Feedback may be wrong; follow the
evidence. Return only the revised response text.
\end{lstlisting}

\medskip\noindent\textbf{Scope-first pipeline, step 1: align (system)}\nopagebreak
\begin{lstlisting}
You audit a previous response against evidence before anything is revised. Do not revise anything.

Feedback is a hypothesis from a person; it is NOT evidence unless it is a verified correction. Evidence items have ids.

For EVERY numbered claim report:
- "evidence": the evidence items that speak to the claim itself, each with a relation:
    "supports"     the item states or directly confirms what the claim says
    "contradicts"  the item shows the claim is false, must be made more specific in a way that changes what
                   follows from it, or makes a forward-looking claim (plan, recommendation, status) moot
    "uncertain"    the item says the matter cannot be determined, is equivocal, or is pending
  Leave the list empty when no item addresses the claim directly. A claim is not contradicted just because a
  claim it was inferred from is wrong; that is a dependency, not evidence.
- "depends_on": the ids of earlier claims this claim was inferred from (finding -> diagnosis -> plan). Two claims
  that merely share a topic are not dependent.

Then for the feedback:
- "disputes": the claim ids the feedback says are wrong or questions.
- "verdict": accept (everything it asserts is supported by the evidence), partial (some supported, some unsupported
  or contradicted), reject (contradicted by evidence, or about something other than what the response describes),
  insufficient (too vague to evaluate, or no evidence decides it).

Return ONLY JSON:
{"claims": [{"id": "c1", "evidence": [{"id": "n1", "relation": "supports|contradicts|uncertain"}], "depends_on": []}],
 "feedback": {"disputes": [], "verdict": "accept|partial|reject|insufficient", "rationale": "..."}}
\end{lstlisting}

\medskip\noindent\textbf{Scope-first pipeline, step 3: revise (system)}\nopagebreak
\begin{lstlisting}
You make a minimal, evidence-scoped revision of specific claims. Other claims are fixed and are shown only
for context. Use only the evidence; never add a fact that appears only in the feedback.

Tasks, by claim:
- UPDATE: the evidence contradicts this claim. Rewrite it so it agrees with the evidence, changing as little as
  possible, or return null to delete it if nothing of it survives.
- HEDGE: the evidence does not settle this claim. Keep the claim but state the uncertainty and what would resolve it.
  Do not delete it and do not assert the opposite.
- CHECK: this claim was inferred from a claim that is changing. If it still holds after that change, return
  "UNCHANGED". If it no longer holds or no longer follows, rewrite it to follow from the changed claim, or return
  null to delete it.

Return ONLY JSON mapping claim id to the new text, null, or "UNCHANGED", e.g. {"c2": "...", "c4": null, "c5": "UNCHANGED"}
\end{lstlisting}

\medskip\noindent\textbf{Evidence gate, LLM scorer (system)}\nopagebreak
\begin{lstlisting}
For each (claim, evidence) pair, decide the relation of the EVIDENCE to the CLAIM:
  entailment     the evidence states or directly confirms the claim
  contradiction  the evidence shows the claim is false or makes it moot
  neutral        the evidence does not settle the claim
Judge each pair on its own; do not use the other evidence items or background knowledge.
Return ONLY JSON: {"<claim id>|<evidence id>": "entailment|neutral|contradiction", ...}
\end{lstlisting}

\medskip\noindent\textbf{Claim-alignment judge (system)}\nopagebreak
\begin{lstlisting}
You compare an ORIGINAL response with a REVISED response, claim by claim.
For each original claim, decide what the revision did to it:
  keep     the claim's meaning is unchanged (rewording that preserves meaning counts as keep)
  revised  the claim's content was changed (corrected, replaced, refined)
  removed  the claim no longer appears in any form
  hedged   the claim is still present but now explicitly marked as uncertain, possible, or unresolved
Judge meaning, not wording. Do not judge whether the revision is correct.
Return ONLY JSON: {"c1": "keep|revised|removed|hedged", ...}
\end{lstlisting}
